\documentclass[11pt,a4paper]{scrartcl}

\usepackage[T1]{fontenc}
\usepackage{lmodern}
\usepackage[ngerman]{babel}
\usepackage{microtype}
\usepackage{amsmath,amssymb}
\usepackage{graphicx}
\usepackage{booktabs}
\usepackage[hidelinks]{hyperref}

\title{{{title}}}
\author{{{author}}}
\date{\today}

\begin{document}
\maketitle

\begin{abstract}
  Eine kurze Zusammenfassung.
\end{abstract}

\section{Einleitung}
\label{sec:einleitung}
Hier beginnt der Text. Quellen werden so zitiert~\cite{knuth1984}.
Abschnitt~\ref{sec:methode} beschreibt das Vorgehen.

\section{Methode}
\label{sec:methode}
\begin{equation}
  a^2 + b^2 = c^2 \label{eq:pythagoras}
\end{equation}

\section{Fazit}
Ergebnisse zusammenfassen.

\bibliographystyle{plaindin}
\bibliography{refs}

\end{document}
